\section{实验设计}

实验设计用于检查所提出的审计工作流能否在可控的数据流通场景中实例化，并使轻量路径能够与全节点确认路径进行对比。实验包含两部分：端到端容器化工作流实验和机制级消融实验。

\subsection{数据集与数据流通负载}

端到端实验使用标准 MNIST 数据集作为可复现实验负载。MNIST 包含 60000 个训练样本和 10000 个测试样本。训练数据被随机划分给十个参与方，用于模拟多方数据流通过程：每个参与方贡献本地更新，而不是直接暴露原始数据。任务运行 100 轮训练。在工作流中，模型更新、审计事件、更新摘要和异常更新记录被视为可审计证据对象。

为测试异常事件处理能力，实验在预设轮次注入恶意更新。实验设置中，第 17、37、57、77 和 97 轮被标记为异常轮次，恶意更新通过将更新方向乘以负因子生成。该注入用于为审计工作流构造可追踪异常事件。

\subsection{容器化区块链审计环境}

系统使用 Docker Compose 部署。容器化环境包括 ZooKeeper、区块链审计节点和 runner。ZooKeeper 维护节点注册表和节点状态。runner 从 ZooKeeper 读取已注册且健康的节点，构造审计载荷，并提交给可用审计节点进行 PBFT 风格投票。每个已提交审计区块记录载荷哈希、Merkle 根、前一区块哈希、投票结果和相关审计元数据。

节点治理实验使用四个审计节点。实验考虑三类治理场景。正常场景中，四个节点均可信且可用。低可信场景中，一个节点仍注册在系统中，但信任分低于阈值，因此被排除出投票集合。故障节点场景中，一个节点已注册且可信，但在确认阶段不可用，剩余节点仍需满足法定人数规则。

\subsection{对比策略}

工作流实验对比四类策略。第一类是无链审计的 ASGD，用作不包含链上审计的训练基线。第二类是 two-factor no-chain 策略，即加入本地异常筛查但不将审计证据提交到链上。第三类是 full PBFT chain 审计，即审计事件由完整可信节点集合确认。第四类是 TrustAuditFlow audit-gate 策略，即常规审计事件使用信任感知门控和链上锚定，而异常或高风险事件被升级处理。

机制级消融进一步将区块链审计路径拆分为四个变体：BLoP 风格全节点 PBFT、带 Merkle 批量锚定的全节点 PBFT、无恢复闭环的委员会确认，以及完整 TrustAuditFlow 路径，即委员会确认、批量锚定、恢复状态记录和回退机制的组合。

\subsection{指标与验证流程}

实验关注审计工作流属性。主要指标包括：链上记录数量、确认延迟、法定人数满足情况、异常事件记录、哈希链连续性、Merkle 根一致性、信任门控下的节点纳入或排除，以及证据冗余。每次运行后，验证器检查已提交区块是否形成连续哈希链、是否满足预期法定人数、异常轮次是否被记录，以及每个场景是否生成可审计链状态。

实验流程如下。首先，准备 MNIST 文件并挂载到 runner 容器。其次，通过 Docker Compose 启动 ZooKeeper 和审计节点。第三，runner 执行选定策略，并在训练轮次中产生审计事件。第四，审计载荷由相应投票集合确认并追加到审计链。随后，验证器检查生成的链文件、节点状态、异常事件记录和摘要指标。该流程为比较全节点审计与轻量化审计门控提供可重复基础。
